\chapter{Phát hiện lỗ hổng ứng dụng web bằng phương pháp Greybox fuzzing}


Chương này trình bày các kiến thức nền tảng cần thiết để hiểu rõ bài toán
và đánh giá giải pháp đề xuất trong khóa luận.
Phần~\ref{sec:greybox-fuzzing} giới thiệu kỹ thuật greybox fuzzing, lịch sử
phát triển và những thách thức đặc thù khi áp dụng lên ứng dụng web PHP.
Phần~\ref{sec:phuzz} mô tả kiến trúc, cơ chế hoạt động và công thức đánh giá
ứng viên của framework Phuzz --- nền tảng mà khóa luận này mở rộng.
Phần~\ref{sec:hooking} mô tả kiến trúc runtime của PHP, cơ chế hook hàm qua
UOPZ và vai trò của phân tích động --- thành phần kỹ thuật cốt lõi để xây
dựng kênh tín hiệu bổ sung cho SSRF.
Phần~\ref{sec:problem-statement} nhận xét về khoảng trống kỹ thuật và phát
biểu bài toán cần giải quyết.

% ============================================================
\section{Greybox Fuzzing}
\label{sec:greybox-fuzzing}
% ============================================================

\subsection{Lịch sử và Bối cảnh Fuzz Testing}

Fuzz testing (hay fuzzing) là kỹ thuật kiểm thử phần mềm tự động được đề
xuất lần đầu bởi Barton Miller vào năm 1988 tại Đại học Wisconsin--Madison,
khi ông nhận thấy rằng các chương trình Unix bị crash khi nhận đầu vào ngẫu
nhiên từ một đường truyền có nhiễu~\cite{miller1990empirical, li2018fuzzing}. Từ xuất phát điểm
giản dị đó, fuzzing đã phát triển thành một trong những kỹ thuật kiểm thử
bảo mật hiệu quả và được nghiên cứu rộng rãi nhất.

Giai đoạn đầu (1988--2010) chứng kiến sự phát triển của các
\emph{blackbox fuzzer} thế hệ đầu: sinh chuỗi ngẫu nhiên hoặc dựa trên
ngữ pháp, không có phản hồi từ chương trình. Mặc dù đơn giản, chúng phát
hiện được nhiều lỗi thực tế trong phần mềm thương mại. Giai đoạn tiếp theo
(2010--nay) đánh dấu bước ngoặt khi American Fuzzy Lop (AFL) của Michal
Zalewski~\cite{zalewski2014afl} phổ biến hóa \emph{coverage-guided fuzzing}:
sử dụng thông tin độ phủ mã thực thi như tín hiệu phản hồi để hướng dẫn
quá trình đột biến. AFL++ \cite{fioraldi2020afl} tiếp tục mở rộng và tổng
hợp nhiều cải tiến nghiên cứu. Song song, Google phát triển
libFuzzer~(in-process) và honggfuzz~\cite{google2023honggfuzz} cho các
kịch bản khác nhau.

Khảo sát của Manes và cộng sự~\cite{manes2021art} phân tích hơn 200 công
trình nghiên cứu về fuzzing và đề xuất một khung thuật ngữ thống nhất, định
nghĩa fuzzing là \textit{``quá trình thực thi PUT (Program Under Test) với
các đầu vào được tạo ra tự động để phát hiện lỗi và/hoặc lỗ hổng bảo mật''}.
Zhu và cộng sự~\cite{zhu2022fuzzing} tổng hợp thêm về roadmap tương lai,
nhấn mạnh sự cần thiết mở rộng fuzzing sang các domain phi truyền thống như
ứng dụng web, hệ thống nhúng và giao thức mạng.

\subsection{Phân loại Fuzzer theo Mức độ Thông tin}

Dựa trên mức độ truy cập thông tin nội bộ của chương trình mục tiêu, fuzzer
được phân thành ba nhóm~\cite{godefroid2007random, li2018fuzzing, zhu2022fuzzing}.

\textbf{Blackbox fuzzer} hoàn toàn không có thông tin về cấu trúc bên trong
ứng dụng. Fuzzer chỉ quan sát được đầu ra bên ngoài (HTTP response, mã trạng
thái, thời gian phản hồi) để đánh giá test case. Ưu điểm là dễ triển khai và
không cần mã nguồn; nhược điểm là khả năng khám phá code path rất hạn chế
trong các ứng dụng có logic phức tạp. Nghiên cứu của Doupé và cộng
sự~\cite{doupe2010johnny} chỉ ra rằng blackbox fuzzer bỏ sót đáng kể các lỗ
hổng trong ứng dụng web, đặc biệt là các lỗ hổng đòi hỏi điều hướng qua
nhiều bước tương tác.

\textbf{Whitebox fuzzer} có đầy đủ thông tin về mã nguồn và cấu trúc chương
trình, cho phép áp dụng các kỹ thuật phân tích tĩnh nặng như thực thi tượng
trưng (symbolic execution)~\cite{godefroid2008automated} hay phân tích vết
lan truyền dữ liệu (taint analysis)~\cite{ganesh2009taint}. Thực thi tượng
trưng biểu diễn các biến đầu vào bằng ký hiệu toán học, xây dựng ràng buộc
từ mỗi nhánh điều kiện, rồi dùng SMT solver để sinh giá trị cụ thể cho phép
chương trình đi theo từng nhánh. Độ chính xác lý thuyết cao là ưu điểm,
nhưng chi phí tính toán lớn (state explosion problem), khó xử lý các thư viện
bên ngoài và mã nguồn động là hạn chế chính khiến kỹ thuật này khó mở rộng
cho ứng dụng web thực tế.

\textbf{Greybox fuzzer} là chiến lược trung gian, đạt được sự đánh đổi tốt
giữa chi phí và hiệu quả: fuzzer thu thập thông tin thực thi (thường là độ
phủ mã) trong khi chạy mà không cần phân tích tĩnh toàn bộ mã nguồn. Thông
tin phủ này trở thành tín hiệu phản hồi để ưu tiên và đột biến các test case
có khả năng khám phá code path mới. AFL++~\cite{fioraldi2020afl} là hiện thực
phổ biến nhất cho ứng dụng nhị phân và là nguồn cảm hứng thiết kế của Phuzz.

\subsection{Coverage-Guided Fuzzing và Vòng lặp Fuzz}
\label{subsec:cgf-loop}

Coverage-guided fuzzing (CGF) là hướng tiếp cận greybox phổ biến
nhất~\cite{li2018fuzzing, manes2021art}. Một test case được coi là
\emph{thú vị} nếu nó kích hoạt ít nhất một code path mới chưa từng được
quan sát trong tập ứng viên hiện tại. Những test case như vậy được giữ lại
trong tập ứng viên (corpus) và ưu tiên đột biến ở các vòng lặp tiếp theo.

Trong fuzzing nhị phân, coverage thường được đo ở đơn vị \emph{edge}
(cạnh trong đồ thị luồng điều khiển, tức là cặp basic block \((A, B)\) nơi
điều khiển chuyển từ $A$ sang $B$). AFL đại diện cho tập edge đã thấy bằng
một \emph{bitmap} $B$ gồm 65536 byte, trong đó byte $B[A \oplus B]$ đếm số
lần edge $(A, B)$ được thực thi trong một test case. Một test case mới được
coi là thú vị nếu bitmap của nó chứa ít nhất một edge chưa có trong bitmap
tổng hợp từ tất cả ứng viên trước đó.

Vòng lặp coverage-guided fuzzing bao gồm sáu bước chính~\cite{manes2021art}:
\begin{enumerate}
    \item \textbf{Khởi tạo}: Nạp tập seed ban đầu vào corpus; thực thi từng
    seed để thiết lập coverage nền $B_0$.
    \item \textbf{Chọn ứng viên}: Áp dụng chiến lược lên lịch seed để chọn
    ứng viên $s$ tiếp theo từ corpus.
    \item \textbf{Phân bổ năng lượng}: Tính số lần đột biến $E(s)$ sẽ dành
    cho $s$ dựa theo chiến lược năng lượng.
    \item \textbf{Đột biến}: Sinh $E(s)$ test case mới $t_1, \ldots, t_{E(s)}$
    từ $s$ bằng các phép đột biến (bit flip, byte insertion, boundary values,\ldots).
    \item \textbf{Thực thi và Đánh giá}: Với mỗi $t_i$: thực thi PUT, thu
    thập bitmap $B_i$; nếu $B_i$ chứa edge mới so với bitmap tổng hợp thì
    thêm $t_i$ vào corpus; nếu $t_i$ kích hoạt hành vi bất thường (crash, lỗi
    bảo mật) thì ghi nhận phát hiện.
    \item Lặp lại từ bước 2 đến khi hết ngân sách thời gian.
\end{enumerate}

Trong fuzzing nhị phân (AFL, AFL++~\cite{fioraldi2020afl, zalewski2014afl},
honggfuzz~\cite{google2023honggfuzz}), coverage được thu thập bằng cách
instrument các basic block tại thời điểm biên dịch (\emph{compile-time
instrumentation}) và chia sẻ bitmap qua vùng nhớ dùng chung (shared memory)
giữa fuzzer và tiến trình PUT.

\begin{figure}[htbp]
    \centering
    % Mô tả hình: Flowchart vòng lặp Coverage-Guided Fuzzing (6 bước, hình tròn/vòng khép kín).
    % Bước 1 (Khởi tạo): hình chữ nhật "Seed pool ban đầu" → thực thi → coverage nền B0.
    % Bước 2 (Chọn ứng viên): hình thoi "Corpus rỗng?" → nếu không: chọn s có score cao nhất.
    % Bước 3 (Phân bổ năng lượng): hình chữ nhật "Tính E(s) theo power schedule".
    % Bước 4 (Đột biến): hình chữ nhật "Sinh E(s) đột biến t1...tE".
    % Bước 5 (Thực thi): hình chữ nhật "Gửi ti đến PUT, thu bitmap Bi".
    % Bước 6 (Đánh giá): hình thoi "Bi có edge mới?" → có: thêm ti vào corpus (quay lại bước 2);
    %   không: hình thoi "ti kích hoạt lỗ hổng?" → có: ghi nhận; không: tiếp tục.
    % Mũi tên vòng lặp quay lại bước 2.
    \includegraphics[width=0.80\textwidth]{figures/fig-cgf-loop}
    \caption{Vòng lặp Coverage-Guided Fuzzing.}
    \label{fig:cgf-loop}
\end{figure}

\subsection{Lên lịch Seed và Phân bổ Năng lượng}
\label{subsec:energy-scheduling}

Chiến lược lên lịch seed (\emph{seed scheduling}) và phân bổ năng lượng
(\emph{energy assignment}) có ảnh hưởng lớn đến tốc độ khám phá code path
mới, nhưng thường bị bỏ qua trong các thiết kế fuzzer đơn giản.

\textbf{AFL gốc} dùng chiến lược round-robin đơn giản: mỗi seed được phân
bổ số lần đột biến bằng nhau, không phân biệt tầm quan trọng. Điều này dẫn
đến lãng phí ngân sách fuzz cho các seed ở vùng code đã được khai thác kỹ.

\textbf{AFLFast}~\cite{bohme2016coverage} đưa ra khái niệm \emph{power schedule}:
mỗi seed $s$ nhận năng lượng $E(s)$ tỷ lệ với \emph{mức độ ít được khám
phá} của vùng code mà $s$ thực thi. Trực giác là: seed thực thi các code path
hiếm (ít test case kích hoạt) nên nhận nhiều năng lượng hơn để fuzzer tập
trung khai thác vùng đó. AFLFast định nghĩa nhiều power schedule khác nhau
(FAST, COE, QUAD, LIN, EXPLOIT, EXPLORE), mỗi cái ưu tiên tiêu chí khác nhau.

Trong Phuzz~\cite{neef2024phuzz}, chiến lược tương tự được áp dụng dựa trên
số đường đi mới và dòng code mới mà mỗi ứng viên khai phá được trong quá
khứ (xem chi tiết trong Phần~\ref{subsec:scoring}).

\subsection{Thách thức khi Áp dụng lên Ứng dụng Web PHP}

Khi chuyển kỹ thuật greybox fuzzing từ ứng dụng nhị phân sang ứng dụng web
PHP, nhiều thách thức đặc thù xuất hiện~\cite{neef2024phuzz,
vanrooij2021webfuzz}.

\textbf{Thách thức về instrumentation.} PHP là ngôn ngữ thông dịch,
không biên dịch trước, nên không thể instrument ở tầng nhị phân như AFL.
Giữa fuzzer và ứng dụng PHP luôn có web server trung gian (Apache, nginx)
tạo lớp gián tiếp. Cơ chế shared memory của AFL không áp dụng được vì các
thành phần chạy trên các container riêng biệt. Phuzz giải quyết bằng cách
dùng PHP extension (Xdebug/PCOV) ghi coverage ra file tmpfs, và fuzzer đọc
lại file sau mỗi request.

\textbf{Thách thức về phát hiện lỗ hổng.} Trong fuzzing nhị phân, crash
của ứng dụng là chỉ số rõ ràng cho hành vi bất thường. Với PHP, interpreter
không crash --- thay vào đó emit exception hoặc error message, và web server
trả về trang lỗi với mã HTTP 5xx. Hơn nữa, ứng dụng có thể bắt và xử lý
lỗi một cách graceful (try/catch), che giấu hoàn toàn tín hiệu. Đặc biệt,
các lớp lỗ hổng như Blind SSRF không để lại bất kỳ dấu vết nào trong HTTP
response --- thách thức này là trọng tâm của khóa luận.

\textbf{Thách thức về sinh test case.} HTTP request có cấu trúc nghiêm ngặt
(method, headers, body encoding); đột biến tùy tiện ở mức byte tạo ra request
không hợp lệ bị web server từ chối trước khi đến PHP. Fuzzer phải thực hiện
đột biến có kiểm soát --- chỉ thay đổi các tham số cụ thể (query string,
POST body, cookie, header) mà không phá vỡ cấu trúc HTTP.

\textbf{Thách thức về trạng thái ứng dụng.} Ứng dụng web thường có trạng
thái phức tạp: session cookie hết hạn, CSRF token thay đổi theo từng request,
database bị ô nhiễm sau các thao tác ghi. Fuzzer cần cơ chế đăng nhập lại tự
động, quản lý session và trong nhiều trường hợp cần reset cơ sở dữ liệu về
trạng thái sạch giữa các lần fuzz để đảm bảo tính tái hiện.

\textbf{Thách thức về môi trường đa tiến trình.} PHP web applications thường
dùng cơ sở dữ liệu, cache và hàng đợi bất đồng bộ. Khi nhiều instance fuzzer
chạy song song, các request có thể can thiệp lẫn nhau qua trạng thái chia
sẻ trong database, dẫn đến hành vi không tái hiện được.

% ============================================================
\section{Framework Phuzz}
\label{sec:phuzz}
% ============================================================

Phuzz~\cite{neef2024phuzz} là framework coverage-guided fuzzing modular,
được thiết kế đặc thù cho ứng dụng web PHP. Được trình bày tại hội nghị
AsiaCCS 2024, Phuzz giải quyết các thách thức nêu trong
Phần~\ref{sec:greybox-fuzzing} thông qua kiến trúc container hóa và cơ chế
instrumentation trong suốt. Đây là nền tảng mà khóa luận này mở rộng để
thêm khả năng phát hiện SSRF.

\subsection{Kiến trúc Tổng thể}

Toàn bộ Phuzz được triển khai trong môi trường Docker~\cite{docker2023}, cho
phép khởi chạy và tái tạo môi trường fuzz một cách nhất quán trên mọi máy.
Hệ thống bao gồm các thành phần chính:

\textbf{Browser và Crawler.} Quá trình bắt đầu bằng việc thu thập các
endpoint cần fuzz. Phuzz hỗ trợ hai phương thức: (1) tự động crawl bằng
headless Chromium qua thư viện Playwright~\cite{microsoft2024playwright},
ghi lại tất cả HTTP request như một HAR file; (2) thủ công bằng cách xuất
HAR file từ DevTools của trình duyệt sau khi tương tác với ứng dụng. Phương
thức thủ công cho phép kiểm soát chính xác phạm vi fuzz, đặc biệt hữu ích
với các API endpoint yêu cầu xác thực phức tạp hoặc multi-step workflow.

\textbf{HARgen và Composegen.} HARgen chuyển đổi HAR file thành file cấu
hình fuzzer ở định dạng JSON. Mỗi tham số (query, body, cookie, header) được
đánh dấu \texttt{fixed} (giữ nguyên, dùng cho session token) hoặc
\texttt{fuzz} (sẽ bị đột biến). Composegen sinh file
\texttt{docker-compose.yml} từ tập cấu hình, cho phép khởi chạy toàn bộ
môi trường chỉ bằng một lệnh. Định dạng cấu hình JSON điển hình:

\begin{verbatim}
{
  "target": "http://web/wp-admin/admin-ajax.php",
  "methods": ["POST"],
  "cookies": {
    "data":  [{"name": "wordpress_logged_in_...",
               "seeds": ["<session>"]}],
    "fixed": [".*"],
    "login": ["wordpress_logged_in_.*"]
  },
  "body_params": {
    "data":   [{"name": "url", "seeds": ["http://example.com"]}],
    "fuzz":   ["url"],
    "weight": 1.0
  }
}
\end{verbatim}

\textbf{Fuzzer.} Thành phần trung tâm, thực hiện vòng lặp fuzz (chi tiết tại
Phần~\ref{subsec:cgf-loop}): đọc cấu hình endpoint, khởi tạo tập ứng viên
ban đầu, chọn ứng viên điểm cao nhất, tính năng lượng, sinh đột biến, gửi
request có header \texttt{X-Fuzzer-Covid: <ID>}, đọc coverage, chạy
VulnChecker và cập nhật điểm ứng viên. Nhiều instance fuzzer có thể chạy
song song, đồng bộ hóa tập ứng viên đã đánh giá qua shared volume để tránh
thực thi trùng lặp.

\textbf{Shared Volume (tmpfs).} Phuzz dùng Docker tmpfs volume làm kênh trao
đổi thông tin giữa các container. PHP-side instrumentation ghi coverage report
vào \texttt{/shared-tmpfs/coverage-reports/<ID>.json}; fuzzer đọc lại để đánh
giá ứng viên. Cơ chế này giải quyết vấn đề thiếu shared memory xuyên container.
Ngoài coverage, tmpfs còn lưu các loại báo cáo khác: hook trigger, OOB log và
kết quả phát hiện lỗ hổng.

\begin{figure}[htbp]
    \centering
    % Mô tả hình: Sơ đồ kiến trúc Phuzz với luồng dữ liệu.
    % Phần trên: Kiểm thử viên → Playwright Browser → HAR file
    %            → HARgen → fuzzer-config.json → Composegen → docker-compose.yml.
    % Docker environment (hình chữ nhật bao quanh 3 container):
    %   - phuzz-fuzzer: "Fuzzer (Python)" — vòng lặp fuzz, VulnChecker, Candidate Pool.
    %   - phuzz-web: "Web Container (PHP 8 + Apache + PCOV)" — ứng dụng mục tiêu,
    %                 auto_prepend_file instrumentation.
    %   - phuzz-db: "Database (MySQL 8)".
    % Mũi tên giữa các container:
    %   fuzzer → [HTTP request + X-Fuzzer-Covid header] → web.
    %   web → [coverage JSON] → shared-tmpfs/coverage-reports/.
    %   fuzzer ← [đọc coverage] ← shared-tmpfs.
    %   web ↔ db (SQL queries).
    % Phần dưới: shared-tmpfs (hình trụ) kết nối fuzzer ↔ web.
    \includegraphics[width=0.88\textwidth]{figures/fig-phuzz-arch}
    \caption{Kiến trúc hệ thống Phuzz.}
    \label{fig:phuzz-arch}
\end{figure}

\subsection{Cơ chế Instrumentation}
\label{subsec:instrumentation}

\textbf{Thu thập coverage.} Phuzz hỗ trợ hai PHP extension:
Xdebug~\cite{rethans2002xdebug} và PCOV~\cite{watkins2023pcov}. Điểm then
chốt là instrumentation hoàn toàn \emph{trong suốt} với ứng dụng mục tiêu
--- không cần chỉnh sửa mã nguồn PHP, cấu hình web server, hay cơ sở dữ
liệu. Phuzz dùng \texttt{auto\_prepend\_file} (cấu hình PHP) để inject
instrumentation code trước khi PHP chạy bất kỳ code ứng dụng nào; và đăng
ký \texttt{register\_shutdown\_function} để đảm bảo coverage luôn được lưu
kể cả khi ứng dụng gọi \texttt{exit()} hoặc bị terminate bởi bộ giới hạn
thời gian.

Mỗi request mang header bổ sung \texttt{X-Fuzzer-Covid: <ID>} (định danh
dạng \texttt{<timestamp>-<UUIDv4>}). PHP-side code đọc header qua
\texttt{\$\_SERVER['HTTP\_X\_FUZZER\_COVID']}, bật coverage collection ngay
đầu request, và lưu report ra \texttt{/shared-tmpfs/coverage-reports/<ID>.json}
khi shutdown. File JSON chứa danh sách file PHP và line number đã được thực thi,
cùng thông tin về lời gọi hàm mà các hook đã ghi lại.

\textbf{So sánh Xdebug và PCOV.}

\begin{itemize}
    \item \textbf{Xdebug}~\cite{rethans2002xdebug}: Extension đa năng (debug,
    profile, trace, coverage); hỗ trợ branch coverage và path coverage theo
    dõi từng nhánh điều kiện. Overhead lớn hơn PCOV nhưng thông tin phong
    phú hơn.
    \item \textbf{PCOV}~\cite{watkins2023pcov}: Extension chuyên dụng cho
    line coverage; overhead thấp hơn đáng kể. Đây là lựa chọn mặc định của
    Phuzz để tối ưu throughput.
\end{itemize}

\textbf{Function hooking với UOPZ.} UOPZ (User Operations for
Zend)~\cite{phpgroup2023uopz} là PHP extension cho phép thao tác hàm ở tầng
interpreter. Phuzz dùng UOPZ để bọc (wrap) các hàm PHP nguy hiểm nhằm bắt
lời gọi hàm, ghi lại đối số và exception:

\begin{verbatim}
uopz_set_return('mysqli_query', function($conn, $query) {
    $result = mysqli_query($conn, $query);
    if ($result === false) {
        $err = mysqli_error($conn);
        phuzz_log_error('sqli', ['query' => $query,
                                 'error' => $err]);
    }
    return $result;
}, true);
\end{verbatim}

Nhờ đó Phuzz có thể bắt lời gọi hàm cùng đối số để xác nhận input từ
attacker đã lan truyền tới điểm nguy hiểm; bắt exception/error do input
không hợp lệ gây ra; và vô hiệu hóa các cơ chế xác thực/phân quyền để
fuzz các code path thông thường không thể tiếp cận.

\subsection{Công thức Đánh giá Ứng viên}
\label{subsec:scoring}

Mỗi ứng viên $c$ trong corpus được gán điểm số để quyết định thứ tự ưu tiên.
Phuzz dùng \texttt{DefaultScoringFormula}, kết hợp hai thành phần:

\begin{itemize}
    \item \textbf{$\Delta\text{paths}$}: Số code path (cặp dòng code liên tiếp
    trong luồng thực thi) mà $c$ khám phá \emph{lần đầu tiên} so với toàn bộ
    corpus hiện tại. Đây là thước đo chính về mức độ ``thú vị'' của $c$.
    \item \textbf{$\Delta\text{lines}$}: Số dòng code mới mà $c$ khám phá, bổ
    trợ cho $\Delta\text{paths}$ trong trường hợp ứng viên cover nhiều dòng
    mới nhưng ít path mới.
\end{itemize}

Điểm tổng hợp được tính theo:
\[
    \text{score}(c) = \alpha \cdot \Delta\text{paths}(c)
                    + \beta  \cdot \Delta\text{lines}(c)
\]
với $\alpha$ và $\beta$ là trọng số cấu hình. Ứng viên có điểm cao nhất được
chọn ở bước 2 của vòng lặp fuzz; khi hai ứng viên bằng điểm, ứng viên mới
hơn được ưu tiên. Năng lượng $E(c)$ --- số lần đột biến --- tỷ lệ tuyến tính
với $\text{score}(c)$, đảm bảo các ứng viên ``thú vị'' nhận được nhiều đột
biến hơn.

\subsection{Phát hiện Lỗ hổng}
\label{subsec:phuzz-vulncheckers}

\textbf{Lỗ hổng phía server.} Phuzz hook các hàm PHP native; khi input từ
fuzzer gây ra lỗi, thông tin debug (tên hàm, đối số, thông báo lỗi) được lưu
trong coverage report và chuyển tới VulnChecker:

\begin{itemize}
    \item \textbf{SQL Injection (SQLi)}: Hook bốn hàm MySQL
    (\texttt{mysqli\_query}, \texttt{mysqli::query}, \texttt{PDO::query},
    \texttt{PDO::exec}); input không hợp lệ tạo ra SQL syntax error được ghi
    nhận trong error log.
    \item \textbf{Command Injection (RCE)}: Hook bốn hàm shell
    (\texttt{shell\_exec}, \texttt{system}, \texttt{passthru}, \texttt{exec});
    kiểm tra thông báo lỗi về lệnh không hợp lệ, hoặc return code khác 0.
    \item \textbf{Path Traversal (PaTr)}: Hook 48 hàm file-related
    (\texttt{file\_get\_contents}, \texttt{fopen}, \texttt{file\_exists},
    v.v.); phát hiện khi path nguy hiểm như \texttt{/etc/passwd} được truyền
    vào hàm file.
    \item \textbf{Insecure Deserialization (IDes)}: Hook \texttt{unserialize()};
    input không hợp lệ tạo ra lỗi PHP object injection.
    \item \textbf{XML External Entity (XXE)}: Hook \texttt{DOMDocument::loadXML}
    trong cấu hình \texttt{LIBXML\_NOENT} cho phép external entity expansion.
\end{itemize}

\textbf{Lỗ hổng phía client.} Hai lớp lỗ hổng client-side được phát hiện
bằng phân tích HTTP response, không cần instrumentation phía server:
Cross-Site Scripting (XSS) theo thuật toán webFuzz~\cite{vanrooij2021webfuzz}
--- chèn marker vào payload và kiểm tra sự xuất hiện của nó trong DOM
response; và Open Redirection (OpRe) --- kiểm tra mã 3xx và giá trị header
\texttt{Location} xem liệu input có xuất hiện trong URL chuyển hướng.

\textbf{ParamBasedVulnChecker.} Để giảm false positive, Phuzz triển khai
\texttt{ParamBasedVulnChecker}: chỉ phát cảnh báo khi ít nhất một tham số
fuzz xuất hiện trong đối số truyền vào hàm bị hook, đảm bảo input từ fuzzer
đã thực sự propagate tới điểm nguy hiểm (sink). Điều này loại bỏ các cảnh
báo giả khi ứng dụng tự tạo ra lỗi SQL không liên quan đến input của fuzzer.

\subsection{Sinh Test Case và Đột biến Tham số}

Phuzz yêu cầu seed ban đầu qua file cấu hình JSON. Ứng viên được chấm điểm
theo số đường đi mới và dòng code mới phía server. Đột biến diễn ra ở mức
byte trên các chuỗi tham số --- không phải toàn bộ HTTP request --- đảm bảo
request luôn hợp lệ. Phuzz định nghĩa kiến trúc \texttt{ParamMutator} plugin:

\begin{itemize}
    \item \textbf{Các mutator byte-level cơ bản}: BitFlipMutator (đảo bit),
    ByteFlipMutator, ArithmeticMutator (tăng/giảm byte value), SpliceMutator
    (ghép hai tham số), DeleteMutator, RepeatMutator, \ldots
    \item \textbf{Mutator ngữ nghĩa đặc thù lỗ hổng}: ProtocolPrefixMutator
    (prefix \texttt{http://}), PathTraversalPayloadParamMutator (chèn
    \texttt{../}), XSSPayloadParamMutator (chèn XSS marker).
    \item \textbf{SSRFMutator} (đề xuất của khóa luận này): chèn 21 biến thể
    URL OOB với các kỹ thuật bypass khác nhau (chi tiết Chương~3).
\end{itemize}

\subsection{Giới hạn của Phuzz và Khoảng trống Nghiên cứu}

Mặc dù đạt hiệu quả cao trong phát hiện các lỗ hổng truyền thống --- 99\%
trong bộ benchmark 87 lỗ hổng đã biết --- bài báo gốc~\cite{neef2024phuzz}
thừa nhận rõ ràng hai giới hạn liên quan trực tiếp đến khóa luận này.

Thứ nhất, Phuzz bỏ sót SSRF trong quá trình fuzzing tự động. Cụ thể,
CVE-2023-6294 (SSRF trong plugin popup-builder cho WordPress) chỉ được phát
hiện thủ công trong bước review sau khi Phuzz tìm thấy lỗ hổng đọc file tùy
ý. Bài báo ghi nhận thẳng thắn: \textit{``The SSRF could have been found by
Phuzz, if \texttt{wp\_remote\_get} was hooked and included in the vulnerability
detection.''} Thứ hai, Phuzz không có kênh phản hồi nào cho các lỗ hổng mà
ứng dụng âm thầm thực hiện hành vi mạng mà không để lại dấu vết trong HTTP
response hay exception. Bài báo liệt kê SSRF là ưu tiên hàng đầu trong phần
Future Work.

Đây chính là khoảng trống mà khóa luận này lấp đầy.

% ============================================================
\section{Phân tích Động và Hook Hàm trong PHP}
\label{sec:hooking}
% ============================================================

\subsection{Kiến trúc Runtime của PHP}
\label{subsec:zend-engine}

PHP được thực thi qua Zend Engine, một trình thông dịch viết bằng C. Hiểu
kiến trúc này là cơ sở để thiết kế cơ chế hooking hiệu quả~\cite{phpgroup2023history}.

Luồng thực thi PHP bao gồm bốn giai đoạn:
\begin{enumerate}
    \item \textbf{Lexing và Parsing}: Mã nguồn PHP được tokenize và parse
    thành Abstract Syntax Tree (AST).
    \item \textbf{Compilation}: AST được biên dịch thành \emph{opcodes}
    (Zend VM bytecode). Opcache có thể cache opcodes để tăng tốc.
    \item \textbf{Execution}: Zend VM thực thi opcode. Mỗi lời gọi hàm tra
    cứu trong \emph{function table} (bảng hash mapping tên hàm thành
    \texttt{zend\_function} struct).
    \item \textbf{Shutdown}: Gọi shutdown functions, giải phóng bộ nhớ.
\end{enumerate}

Điểm quan trọng là \emph{function table} có thể được thao tác ở runtime bởi
các extension, cho phép thay thế hoặc bọc bất kỳ hàm PHP built-in nào. Đây
là cơ chế mà UOPZ khai thác.

\textbf{Cơ chế auto\_prepend\_file.} PHP cung cấp directive
\texttt{auto\_prepend\_file} trong \texttt{php.ini}: một file PHP được include
tự động trước khi thực thi bất kỳ script PHP nào. Phuzz dùng cơ chế này để
inject instrumentation code (coverage collection và function hooks) một cách
trong suốt --- ứng dụng mục tiêu không biết về sự tồn tại của code bổ sung
này.

\subsection{Extension UOPZ --- Cơ chế và API}
\label{subsec:uopz}

UOPZ (User Operations for Zend Engine)~\cite{phpgroup2023uopz} là PHP extension
cho phép thao tác với hàm và lớp PHP tại tầng Zend Engine ở runtime, không
cần sửa mã nguồn. Ba API chính phục vụ các mục đích khác nhau:

\begin{itemize}
    \item \textbf{uopz\_set\_hook(\$function, \$hook)}: Đặt hook được gọi
    \emph{trước} hàm gốc, với cùng tham số. Dùng để capture đối số trước khi
    hàm thực thi; hàm gốc vẫn chạy bình thường sau đó. Đây là API dùng trong
    khóa luận này để hook các hàm mạng SSRF.
    \item \textbf{uopz\_set\_return(\$function, \$value, \$execute = false)}:
    Thay thế return value của hàm. Khi \texttt{\$execute = true}, wrap hàm
    bằng closure và gọi hàm gốc bên trong. Phuzz dùng API này để bắt exception.
    \item \textbf{uopz\_redefine(\$constant, \$value)}: Redefine hằng số tại
    runtime.
    \item \textbf{uopz\_undefine(\$function)}: Vô hiệu hóa hoàn toàn một hàm.
    Phuzz dùng để disable các hàm xác thực/phân quyền.
\end{itemize}

Ví dụ về \texttt{uopz\_set\_hook} để ghi nhận lời gọi \texttt{file\_get\_contents}:

\begin{verbatim}
uopz_set_hook('file_get_contents', function($filename) {
    // Chạy trước file_get_contents() thực sự
    if (filter_var($filename, FILTER_VALIDATE_URL)) {
        phuzz_log_ssrf_hook([
            'func'     => 'file_get_contents',
            'url'      => $filename,
            'covID'    => $_SERVER['HTTP_X_FUZZER_COVID'] ?? '',
        ]);
    }
});
\end{verbatim}

Ưu điểm quan trọng nhất của UOPZ là hook hoạt động \emph{ở tầng interpreter},
nghĩa là áp dụng cho mọi lời gọi từ bất kỳ tầng abstraction nào: dù ứng dụng
gọi hàm trực tiếp, qua framework, hay qua nhiều lớp wrapper, hook vẫn được
kích hoạt.

\textbf{So sánh với các cơ chế hooking khác.}
\begin{itemize}
    \item \textbf{Xdebug trace}: Ghi lại mọi lời gọi hàm theo kiểu trace,
    nhưng không cho phép can thiệp vào luồng thực thi hay sửa đổi đối số.
    Overhead rất lớn.
    \item \textbf{APM agents (Datadog, New Relic)}: Dùng cơ chế tương tự để
    instrument function calls, nhưng được thiết kế cho performance monitoring,
    không phải security testing.
    \item \textbf{Monkey patching (override trong PHP user-space)}: Không thể
    override PHP built-in functions bằng cách đơn giản như trong Python ---
    UOPZ là cần thiết cho mục đích này.
\end{itemize}

\subsection{Phân loại Hàm PHP Có Nguy cơ SSRF}

Trong PHP, SSRF không chỉ phát sinh từ một hàm duy nhất mà từ nhiều nhóm
hàm khác nhau, mỗi nhóm có đặc điểm riêng về cách thực hiện kết nối
mạng. Hiểu rõ các nhóm này là cơ sở để thiết kế cơ chế hook toàn diện.

\textbf{Nhóm Network/HTTP.} Các hàm như \texttt{file\_get\_contents},
\texttt{fopen}, \texttt{readfile}, \texttt{file} là các hàm đa năng trong
PHP --- khi nhận một URL thay cho đường dẫn file cục bộ, chúng âm thầm thực
hiện HTTP/FTP request thông qua stream wrapper. Đây là nhóm phổ biến nhất vì
lập trình viên thường dùng chúng để đọc dữ liệu từ remote API mà không nhận
ra nguy cơ. Các hàm \texttt{get\_headers}, \texttt{get\_meta\_tags},
\texttt{getimagesize} cũng thực hiện HTTP request ngầm khi nhận URL.

\textbf{Nhóm cURL.} Thư viện cURL (\texttt{curl\_exec},
\texttt{curl\_multi\_exec}) cung cấp một client HTTP đầy đủ tính năng. Điểm
đặc biệt về kỹ thuật: URL được thiết lập qua \texttt{curl\_setopt} ở bước
cấu hình (\texttt{CURLOPT\_URL}), không phải qua \texttt{curl\_exec} ở bước
thực thi --- do đó cần hook cả hai để không bỏ sót URL được set từ trước.

\textbf{Nhóm Sockets.} \texttt{fsockopen}, \texttt{pfsockopen} và
\texttt{stream\_socket\_client} mở kết nối TCP/UDP thô trực tiếp tới
hostname:port tùy ý, bỏ qua mọi lớp HTTP abstraction. Đây thường là điểm
sink trong các thư viện HTTP bậc thấp và protocol client (SMTP, Redis, \ldots).

\textbf{Nhóm Shell execution.} Ngoài SSRF thuần túy, một kịch bản nguy hiểm
khác là Command Injection dẫn tới SSRF: input của người dùng được truyền vào
các hàm thực thi shell (\texttt{shell\_exec}, \texttt{system}, \texttt{exec},
\texttt{passthru}, \texttt{popen}) và chứa lệnh như \texttt{curl} hay
\texttt{wget} trỏ tới địa chỉ tùy ý. Việc hook nhóm này cho phép phát hiện
đồng thời cả hai lớp lỗ hổng: Command Injection và SSRF qua command.

\textbf{Nhóm CMS wrappers.} Các framework như WordPress cung cấp các hàm
HTTP wrapper cấp cao (\texttt{wp\_remote\_get}, \texttt{wp\_remote\_post},
\texttt{wp\_remote\_request}, \texttt{wp\_remote\_head}, \texttt{wp\_safe\_remote\_get},
\texttt{wp\_safe\_remote\_post}) nằm trên cùng của thư viện WP HTTP. Nhiều
plugin WordPress dùng các hàm này thay cho cURL hay \texttt{file\_get\_contents}
trực tiếp; nếu chỉ hook tầng thấp sẽ bỏ sót các lời gọi thông qua wrapper.

\textbf{Nhóm Stream wrappers.} PHP hỗ trợ \texttt{stream\_socket\_client},
\texttt{stream\_context\_create} cùng với \texttt{fopen} để tạo kết nối mạng
bậc thấp. Điểm đặc biệt là stream wrapper có thể được đăng ký tùy chỉnh bởi
các extension (ví dụ: \texttt{s3://} cho Amazon S3), tạo ra các sink SSRF
không có trong danh sách hook mặc định.

Một thiết kế hook SSRF hoàn chỉnh cần bao phủ cả năm nhóm trên. Chi tiết
cài đặt cụ thể được trình bày trong Chương~3.

\subsection{Kiến trúc Hai Tầng Tín hiệu cho SSRF}
\label{subsec:two-layer}

Giải pháp đề xuất trong khóa luận này kết hợp hai tầng tín hiệu bổ sung cho
nhau, giải quyết các giới hạn của từng tầng đơn lẻ:

\textbf{Tầng 1 --- PHP Hook (In-band, đồng bộ).} Hook các hàm mạng PHP bằng
UOPZ để ghi nhận lời gọi hàm ngay trong request handler. Khi hàm mạng được
gọi với một URL, hook ghi thông tin (tên hàm, URL đích, coverage ID) vào
\texttt{/shared-tmpfs/ssrf-hooks/<covID>.json}. Tín hiệu nhanh và trực tiếp,
cung cấp ngữ cảnh phong phú (hàm nào, URL nào), nhưng chỉ kích hoạt khi
ứng dụng thực sự gọi hàm mạng với input của fuzzer và hook được kích hoạt
trong cùng request handler.

\textbf{Tầng 2 --- OOB Callback (Out-of-band, bất đồng bộ).} Lắng nghe kết
nối từ ứng dụng mục tiêu trên OOB server. Khi ứng dụng thực hiện HTTP request
tới URL OOB (\texttt{http://oob-server/<token>}), OOB server ghi log vào
\texttt{/shared-tmpfs/oob-logs/<token>.json}. Đây là bằng chứng mạng không
thể phủ nhận: nếu OOB server nhận kết nối mang token của fuzzer, ứng dụng đã
thực sự thực hiện network request tới địa chỉ ngoài. Tín hiệu này đặc biệt
quan trọng cho Blind SSRF với xử lý bất đồng bộ.

Hai tầng này được kết hợp trong \texttt{SSRFVulnCheck} theo logic \textsc{AND}:
một ứng viên chỉ được báo cáo là SSRF khi có \emph{cả hai} tín hiệu cùng
hiện diện --- hook ghi nhận lời gọi hàm mạng, \emph{và} OOB server nhận được
callback. Điều này đảm bảo:
\begin{itemize}
    \item Tầng 1 cung cấp phát hiện nhanh và ngữ cảnh (URL đích, hàm được
    gọi, line số trong PHP).
    \item Tầng 2 là xác nhận cuối cùng không thể tranh cãi: bằng chứng mạng
    thực tế.
    \item Logic AND loại bỏ false positive từ cả hai phía: hook có thể bị
    kích hoạt bởi request hợp lệ, nhưng OOB callback chỉ xảy ra khi ứng dụng
    thực sự kết nối đến OOB server với đúng token.
\end{itemize}

% ============================================================
\section{Nhận xét về Bài toán cần Giải quyết}
\label{sec:problem-statement}
% ============================================================

Dựa trên các kiến thức nền tảng được xây dựng trong chương này, có thể
phát biểu bài toán nghiên cứu một cách chính xác và đầy đủ.

\subsection{Tổng hợp Khoảng trống Kỹ thuật}

Greybox fuzzing với coverage guidance (Phần~\ref{sec:greybox-fuzzing}) là
phương pháp phù hợp nhất để kiểm thử bảo mật ứng dụng web PHP tự động:
nó không yêu cầu mã nguồn, mở rộng tốt, và có thể tích hợp tín hiệu phản
hồi phong phú. Phuzz~\cite{neef2024phuzz} (Phần~\ref{sec:phuzz}) là hiện
thực hóa tốt nhất của hướng tiếp cận này, đạt 99\% phát hiện trên bộ
benchmark 87 lỗ hổng truyền thống.

Tuy nhiên, kiến trúc hiện tại của Phuzz --- và của mọi greybox fuzzer web
hiện tại --- dựa trên giả định cơ bản rằng tín hiệu phản hồi đến từ hai
kênh: HTTP response và code coverage. Cả hai kênh đều không thể quan sát
hành vi mạng ẩn của ứng dụng. Bảng~\ref{tab:signal-gap} tổng hợp khoảng
trống này:

\begin{table}[h]
\centering
\caption{Tín hiệu phản hồi sẵn có và khả năng phát hiện từng loại lỗ hổng}
\label{tab:signal-gap}
\begin{tabular}{|l|c|c|c|c|}
\hline
\textbf{Tín hiệu phản hồi} & \textbf{SQLi} & \textbf{XSS} & \textbf{SSRF} & \textbf{Blind SSRF} \\
\hline
Nội dung HTTP response      & \checkmark & \checkmark & \checkmark & \ding{55} \\
HTTP status code bất thường & \checkmark & \ding{55}  & Có thể     & \ding{55} \\
Exception / thông báo lỗi   & \checkmark & \ding{55}  & Có thể     & \ding{55} \\
Độ phủ mã (code coverage)   & \checkmark & \checkmark & \checkmark & \ding{55} \\
\textbf{OOB network callback}& \ding{55}  & \ding{55}  & \checkmark & \checkmark \\
\hline
\end{tabular}
\end{table}

Cơ chế hook hàm PHP qua UOPZ (Phần~\ref{sec:hooking}) cung cấp khả năng
bắt lời gọi hàm mạng ở tầng interpreter --- nhưng bản thân hook chỉ là
tín hiệu in-band, không thể xác nhận ứng dụng đã \emph{thực sự} thực hiện
kết nối ra ngoài.

\subsection{Phát biểu Bài toán}

Bài toán nghiên cứu của khóa luận này là:

\begin{quote}
\textit{Thiết kế và cài đặt một hệ thống phát hiện Blind SSRF hoàn toàn
tự động, tích hợp vào vòng lặp greybox fuzzing PHP của Phuzz, sao cho:
(1)~mọi tham số HTTP có thể là SSRF sink đều được kiểm thử với payload
bao phủ các kỹ thuật bypass phổ biến; (2)~bằng chứng lỗ hổng được xác
nhận qua hai kênh tín hiệu độc lập (in-band hook và out-of-band callback)
theo logic AND, đảm bảo zero false positive; và (3)~toàn bộ quá trình từ
sinh payload đến xác nhận lỗ hổng diễn ra tự động mà không cần can thiệp
thủ công.}
\end{quote}

\subsection{Yêu cầu Kỹ thuật}

Bài toán đặt ra ba yêu cầu kỹ thuật cụ thể cần giải quyết:

\begin{enumerate}
    \item \textbf{Kênh tín hiệu thứ ba (OOB).} Cần một máy chủ lắng nghe
    độc lập có thể nhận kết nối từ ứng dụng mục tiêu và tương quan chính
    xác mỗi callback với ứng viên fuzz đã kích hoạt nó --- ngay cả trong
    môi trường song song với hàng nghìn request đồng thời.
    \item \textbf{Bộ đột biến chuyên biệt.} Cần payload SSRF bao phủ đầy
    đủ các kỹ thuật bypass (mã hóa IP, DNS redirect, userinfo trick, \ldots)
    để không bỏ sót các lỗ hổng chỉ phát hiện được qua một biến thể cụ thể.
    \item \textbf{Tích hợp không xâm lấn.} Giải pháp phải tương thích với
    kiến trúc plugin của Phuzz, không phá vỡ các chức năng phát hiện lỗ
    hổng gốc, và có thể triển khai trên bất kỳ ứng dụng PHP mục tiêu nào
    mà không cần sửa đổi mã nguồn ứng dụng.
\end{enumerate}

Chương~3 trình bày giải pháp đáp ứng đầy đủ cả ba yêu cầu trên.

\section{Tổng kết Chương}

Chương này đã xây dựng nền tảng lý thuyết cho toàn bộ khóa luận qua ba
mảng kiến thức liên kết chặt chẽ.

Greybox fuzzing với coverage guidance là chiến lược hiệu quả để khám phá
lỗ hổng trong ứng dụng web PHP, nhưng yêu cầu giải quyết các thách thức
đặc thù: instrumentation qua HTTP, quản lý trạng thái ứng dụng và phát hiện
lỗ hổng không có crash signal. Phuzz là framework hiện thực hóa chiến lược
đó, đạt hiệu quả cao trên nhiều lớp lỗ hổng truyền thống và được thiết kế
modular để mở rộng. Tuy nhiên, kiến trúc hiện tại của Phuzz không có khả
năng phát hiện Blind SSRF vì cả hai kênh phản hồi truyền thống (HTTP response
và coverage) đều không quan sát được hành vi mạng ẩn của ứng dụng.

Cơ chế hook hàm mạng PHP qua UOPZ cung cấp tín hiệu in-band tức thì khi
hàm mạng được gọi. Tuy nhiên hook đơn lẻ chưa đủ --- cần một kênh tín hiệu
ngoài băng (OOB) độc lập để xác nhận ứng dụng đã thực sự kết nối ra ngoài.

Hai tầng tín hiệu (PHP hook AND OOB callback) cùng tạo thành nền tảng kỹ
thuật cho thiết kế và cài đặt được trình bày trong Chương~3, nơi lỗ hổng
SSRF và kỹ thuật OOB được phân tích chi tiết trước khi trình bày cài đặt
hệ thống.
